\documentclass[11pt,a4paper]{article}

\usepackage[utf8]{inputenc}
\usepackage[margin=1in]{geometry}
\usepackage{amsmath,amssymb}
\usepackage{fancyhdr}
\usepackage{xcolor}

% Header
\pagestyle{fancy}
\fancyhf{}
\lhead{\textbf{Modern Computer Vision} --- Technion, Spring 2026}
\rhead{HW 2 --- Theory}
\renewcommand{\headrulewidth}{0.4pt}

\newcommand{\R}{\mathbb{R}}

\begin{document}

\begin{center}
  {\LARGE\bfseries Homework 2 --- Theory}\\[6pt]
  \textcolor{gray}{Due: May 10, 2026}
\end{center}

\vspace{12pt}

\noindent
Consider a layer in a neural network whose input is an invertible matrix $A \in \R^{n \times n}$ and whose output is $A^{-1}$.
What is the backward of this layer $\bigl(\frac{\partial \mathcal{L}}{\partial A}\bigr)$?

\medskip
\noindent
\textcolor{gray}{\small
Guidance: first, express the Jacobian $\frac{\partial A^{-1}}{\partial A}$
(think about what type of object it is).
Then use the chain rule to obtain $\frac{\partial \mathcal{L}}{\partial A}$.
You may work with tensor operations or with index notation (e.g.\ $A_{ij}$)---whichever you prefer.}

\vfill
\begin{center}
{\small\textcolor{gray}{Modern Computer Vision --- Technion --- Spring 2026 --- Assaf Shocher}}
\end{center}

\end{document}
